\documentclass[11pt]{article}

\usepackage[margin=1in]{geometry}
\usepackage{amsmath, amssymb, amsthm}
\usepackage{mathtools}
\usepackage{bm}

\DeclareMathOperator*{\clip}{clip}
\DeclareMathOperator{\sign}{sgn}
\newcommand{\Ehat}{\widehat{\mathbb{E}}}
\newcommand{\smean}{s^{\text{MEAN}}}

\newtheorem{theorem}{Theorem}
\newtheorem{definition}{Definition}
\newtheorem{lemma}{Lemma}
\newtheorem{proposition}{Proposition}
\newtheorem{corollary}{Corollary}

\title{Canonical Hierarchical Decomposition of the BSPO Ratio Statistics \\
\large --- A Derivation Independent of GRPO's Zero-Group-Mean Property}
\author{}
\date{}

\begin{document}
\maketitle

\begin{abstract}
We drop the assumption that GRPO advantages satisfy
$\sum_{\tau\in g}\hat A_\tau = 0$ and ask: among the candidate
normalisations A1, A2, A3 of $s_\tau^\pm$, $s_g^\pm$, $s_{dom}^\pm$, which
one produces the \emph{canonical} hierarchical decomposition that retains
the recovery/decoupling property when the cross-level subtraction
$s_g - s_{dom}$ (or $s_\tau - s_g$) is used in contexts outside the
$\Ehat\cdot\hat A_\tau$ linear form --- e.g.\ penalty terms, or novel
variants? We show that A3 (mean at every level) is uniquely singled out by
an ANOVA-style decomposition criterion: it is the only one of the three in
which the empirical $s_g$ equals the empirical conditional mean of
$s_\tau$ given $g$, and similarly for $s_{dom}$. A2 fails this uniqueness
even though it passes the narrower recovery test of
\texttt{bisim\_eqs\_recovery.tex}. The argument does not use
$\sum_{\tau\in g}\hat A_\tau=0$ at any point.
\end{abstract}

% =============================================================================
\section{Setup}

As in \texttt{bisim\_eqs\_recovery.tex}, denote the three candidate
normalisations of the one-sided ratio statistics by A1, A2, A3. The
trajectory-level definition agrees on A2 and A3
($s_\tau^\pm = |\tau|^{-1}\sum_t\max(\pm(r_t-1),0)$), and differs on A1
(no division). The group- and domain-level definitions are the point of
interest:
\begin{align*}
\text{A2:}\quad s_g^\pm &= \sum_{\tau\in g}\max(\pm s_\tau,0),
  & s_{dom}^\pm &= \sum_{g\in dom}\max(\pm s_g,0) \\
\text{A3:}\quad s_g^\pm &= \tfrac{1}{|g|}\sum_{\tau\in g}\max(\pm s_\tau,0),
  & s_{dom}^\pm &= \tfrac{1}{|dom|}\sum_{g\in dom}\max(\pm s_g,0).
\end{align*}

\textbf{New assumption.} We \emph{drop} the GRPO identity
$\sum_{\tau\in g}\hat A_\tau = 0$. All cancellations in this note come
purely from the algebraic aggregation structure of the advantage
($\hat A_g = |g|^{-1}\sum_\tau\hat A_\tau$;
$\hat A_{dom} = |dom|^{-1}\sum_g\hat A_g$ under balanced batches, or the
$T_g/T_{dom}$-weighted form of \texttt{bisim\_eqs\_multidomain.tex}~\S7
in general) and from the normalisation of $s_z^\pm$.

% =============================================================================
\section{Two contexts in which $s_g - s_{dom}$ may appear}

\subsection*{Context I --- multiplied by a hierarchical advantage}

This is the Variant~2 full form of \texttt{bisim\_eqs\_multidomain.tex}~\S8:
\[
\hat J_2^{\text{full}}
  = \Ehat_\tau\bigl[\Delta_{dom}\hat A_{dom}
       + (\Delta_g-\Delta_{dom})\hat A_g
       + (\Delta_\tau-\Delta_g)\hat A_\tau\bigr].
\]
At $\delta\to\infty$, $\Delta_z\to s_z$.

\begin{lemma}[Aggregation identities, no GRPO needed]
\label{lem:agg}
For any function $s_z$ that is constant on the $z$-block of the hierarchy
(trajectory, group, or domain), and with $\hat A_g=|g|^{-1}\sum_\tau\hat A_\tau$
and balanced batches,
\[
\Ehat_\tau[s_{g(\tau)}\hat A_\tau]
  = \Ehat_\tau[s_{g(\tau)}\hat A_{g(\tau)}],
\qquad
\Ehat_\tau[s_{dom(\tau)}\hat A_g]
  = \Ehat_\tau[s_{dom(\tau)}\hat A_{dom(\tau)}].
\]
\end{lemma}

\begin{proof}
$\Ehat_\tau[s_{g(\tau)}\hat A_\tau] = B^{-1}\sum_\tau s_{g(\tau)}\hat A_\tau
= B^{-1}\sum_g s_g\sum_{\tau\in g}\hat A_\tau = B^{-1}\sum_g s_g\cdot|g|\hat A_g
= B^{-1}\sum_g\sum_{\tau\in g}s_g\hat A_g = \Ehat_\tau[s_{g(\tau)}\hat A_{g(\tau)}]$.
Same for the other identity.
\end{proof}

\begin{proposition}[Context I recovery holds without GRPO]
\label{prop:ctx1}
For both A2 and A3, at $\delta\to\infty$,
\[
\hat J_2^{\text{full},(\infty)}
  = \Ehat_\tau[s_\tau\hat A_\tau]
  = \hat J_1^{(\infty)}.
\]
\end{proposition}

\begin{proof}
Expanding and applying Lemma~\ref{lem:agg} twice:
\begin{align*}
\hat J_2^{\text{full},(\infty)}
  &= \Ehat[s_{dom}\hat A_{dom}]
   + \Ehat[s_g\hat A_g] - \Ehat[s_{dom}\hat A_g]
   + \Ehat[s_\tau\hat A_\tau] - \Ehat[s_g\hat A_\tau] \\
  &= \Ehat[s_{dom}\hat A_{dom}]
   + \Ehat[s_g\hat A_g] - \Ehat[s_{dom}\hat A_{dom}]
   + \Ehat[s_\tau\hat A_\tau] - \Ehat[s_g\hat A_g] \\
  &= \Ehat[s_\tau\hat A_\tau].
\end{align*}
No use of $\sum\hat A_\tau|g = 0$.
\end{proof}

\textbf{Consequence for Context I.} The full-hierarchy form of Variant 2
recovers Variant 1 under both A2 and A3, purely by aggregation algebra,
independent of GRPO. So in Context I the recovery property \emph{alone}
cannot discriminate A2 from A3 (this matches the conclusion of
\texttt{bisim\_eqs\_recovery.tex}~\S6).

\subsection*{Context II --- the subtraction appears without $\hat A$}

Everywhere else in the BSPO objective collection the higher-level $s_z$
appears \emph{unmultiplied} by an advantage. The three places in the
current spec that matter:

\begin{enumerate}
\item \emph{Variant 4--5 penalty terms}:
\[
-\lambda_z \max\!\Bigl(0,\;\tfrac{\min(s_z^+,s_z^-)}{\delta_z}-1\Bigr),
\qquad z\in\{\tau,g,dom\}.
\]
The $\min$-based clipping is compared to $\delta_z$ directly; there is no
$\hat A$ in sight. The scale of $s_z$ enters through the ratio
$\min(s_z^+,s_z^-)/\delta_z$.

\item \emph{Variant 2 simplified form}:
$\Ehat[(\Delta_\tau-\Delta_g)\hat A_\tau]$. At $\delta\to\infty$ this is
$\Ehat[s_\tau\hat A_\tau] - \Ehat[s_g\hat A_\tau]$. Without GRPO
zero-mean, $\Ehat[s_g\hat A_\tau] = \Ehat[s_g\hat A_g] \neq 0$ in general
(Lemma~\ref{lem:agg}). Recovery to $\hat J_1$ \emph{fails in the simplified
form} whenever the per-group mean of $\hat A$ is nonzero; the size of the
failure is governed by $|s_g|\cdot|\Ehat[s_g\hat A_g]|$.

\item \emph{Hypothetical novel variants} where the cross-level residual
itself appears as a standalone regulariser, e.g.\
$-\lambda\max(0, |s_g - s_{dom}|/\delta - 1)$ or an ANOVA-style term
$\sum_z \|s_z - s_{\text{parent}(z)}\|^2$. Such terms do not involve
$\hat A$ at all.
\end{enumerate}

In every item of Context II, the magnitude of $s_g$ (resp.\ $s_{dom}$) is
not absorbed into a multiplication-by-advantage, and therefore the
normalisation choice \emph{does} change the numerical behaviour of the
objective.

\begin{proposition}[Context II separates A2 from A3]
\label{prop:ctx2}
Under A2, $|s_g| = |g|\cdot|\smean_g|$ and $|s_{dom}| = |dom|\cdot|g|\cdot|\smean_{dom}|$
where $\smean_g = |g|^{-1}\sum_{\tau\in g}s_\tau$ and
$\smean_{dom} = (|dom|\cdot|g|)^{-1}\sum_{\tau\in dom}s_\tau$. Under A3,
$|s_g| = |\smean_g|$ and $|s_{dom}| = |\smean_{dom}|$. Consequently, for
any of the Context-II constructions above:
\begin{itemize}
\item A2 forces each level's hyperparameter ($\delta_z$ or $\lambda_z$) to
carry the corresponding level-size factor ($|g|$, $|g|\cdot|dom|$). The
hyperparameter is no longer a property of the policy-optimisation problem
but a property of the \emph{batch composition}.
\item A3 places all three levels on the same O($\varepsilon$) scale, so a
single fundamental hyperparameter value applies uniformly, invariant to
batch composition.
\end{itemize}
\end{proposition}

This is the \emph{weak form} of the distinction. The user's objection
``although the scale can be absorbed into $\lambda$'' is correct in a
literal sense: A2 is not \emph{mathematically} wrong, one can always
rescale hyperparameters. But the rescaling required is not a free choice;
it is $(|g|, |dom|\cdot|g|)$-dependent, and therefore not canonical. The
strong form of the distinction below makes the canonical argument.

% =============================================================================
\section{Canonical decomposition --- the unique characterisation of A3}

Let $P$ denote the joint sampling distribution over trajectories, groups,
and domains. Define the population-level quantities
\[
\mu_g := \mathbb{E}_{\tau\sim P(\cdot\mid g)}[s_\tau], \qquad
\mu_{dom} := \mathbb{E}_{g\sim P(\cdot\mid dom)}[\mu_g].
\]

\begin{definition}[Canonical decomposition]
\label{def:canonical}
Given empirical estimators $(\widehat s_\tau,\widehat s_g,\widehat s_{dom})$,
call the collection \emph{canonical} if all of the following hold for
every batch:
\begin{enumerate}
\item[(C1)] $\widehat s_\tau$ is a consistent estimator of
$\mathbb{E}_{t\sim\text{Unif}(\tau)}[r_t - 1]$, the within-trajectory
population ratio drift.
\item[(C2)] $\widehat s_g$ is the empirical conditional mean of
$\widehat s_\tau$ over $\tau\in g$: $\widehat s_g = |g|^{-1}\sum_{\tau\in g}\widehat s_\tau$.
\item[(C3)] $\widehat s_{dom}$ is the empirical conditional mean of
$\widehat s_g$ over $g\in dom$: $\widehat s_{dom} = |dom|^{-1}\sum_{g\in dom}\widehat s_g$.
\end{enumerate}
\end{definition}

\begin{theorem}[A3 is the unique canonical decomposition among A1--A3]
\label{thm:unique}
Of the three candidates,
\begin{itemize}
\item A1 violates (C1): $\widehat s_\tau^{\text{A1}} = |\tau|\cdot\smean_\tau$,
which depends on the trajectory length and is not a consistent estimator
of $\mathbb{E}[r_t-1]$ as $|\tau|\to\infty$.
\item A2 satisfies (C1) but violates (C2) and (C3):
$\widehat s_g^{\text{A2}} = \sum_{\tau\in g}\max(\pm s_\tau,0) \neq |g|^{-1}\sum_\tau s_\tau$
whenever $|g|\neq 1$.
\item A3 satisfies (C1), (C2), and (C3) by construction.
\end{itemize}
\end{theorem}

\begin{proof}
Direct from the definitions. For (C2) under A3: using the identity
$\max(x,0)-\max(-x,0)=x$,
\begin{align*}
\widehat s_g^{\text{A3}} &= \widehat s_g^{+,\text{A3}} - \widehat s_g^{-,\text{A3}}
= \tfrac{1}{|g|}\sum_{\tau\in g}\max(s_\tau,0) - \tfrac{1}{|g|}\sum_{\tau\in g}\max(-s_\tau,0) \\
&= \tfrac{1}{|g|}\sum_{\tau\in g}\bigl[\max(s_\tau,0) - \max(-s_\tau,0)\bigr]
= \tfrac{1}{|g|}\sum_{\tau\in g}s_\tau. \qed
\end{align*}
The same identity gives (C3).
\end{proof}

% =============================================================================
\section{What (C1--C3) buy us that A2 does not}

\subsection*{4.1 Law of total expectation at the empirical level}

Under A3, (C2) and (C3) give the empirical law-of-total-expectation:
\[
\widehat s_{dom}
= |dom|^{-1}\sum_{g\in dom}\widehat s_g
= |dom|^{-1}\sum_{g\in dom}|g|^{-1}\sum_{\tau\in g}s_\tau
= \underbrace{\bigl(|dom|^{-1}|g|^{-1}\bigr)}_{\text{weight}}\sum_{\tau\in dom}s_\tau,
\]
which, for balanced groups, coincides with the unweighted domain mean
$(|dom|\cdot|g|)^{-1}\sum_{\tau\in dom}s_\tau$. Under A2 no such tower
identity holds: $\widehat s_{dom}^{\text{A2}} = \sum_g\sum_\tau |s_\tau| \neq (\text{constant})\sum_\tau s_\tau$.

\subsection*{4.2 ANOVA decomposition}

\begin{proposition}
\label{prop:anova}
Under A3, for any $\tau \in g \in dom$,
\[
s_\tau \;=\; \widehat s_{dom} \;+\; (\widehat s_g - \widehat s_{dom}) \;+\; (s_\tau - \widehat s_g),
\]
and the three terms on the right form an orthogonal decomposition in the
sense of Analysis of Variance:
\begin{enumerate}
\item $\Ehat_{\tau\in g}[s_\tau - \widehat s_g] = 0$ for every $g$ (within-group).
\item $\Ehat_{g\in dom}[\widehat s_g - \widehat s_{dom}] = 0$ for every $dom$ (between-group within-domain).
\item $\widehat s_{dom}$ is the unconditional mean at the domain level.
\end{enumerate}
\end{proposition}

\begin{proof}
Identities (1) and (2) follow directly from (C2) and (C3). The decomposition
in the first line is trivially the telescoping sum. Orthogonality in the
ANOVA sense follows from the fact that each component is a residual
against the next level's mean.
\end{proof}

\begin{proposition}[A2 does not admit the ANOVA decomposition]
Under A2, $\Ehat_{\tau\in g}[s_\tau - \widehat s_g^{\text{A2}}] =
\widehat s_g^{\text{A3}} - \widehat s_g^{\text{A2}}
= (1-|g|)\cdot\widehat s_g^{\text{A3}} \neq 0$ whenever $|g|>1$ and
$\widehat s_g^{\text{A3}}\neq 0$.
\end{proposition}

\subsection*{4.3 Consistency of estimators at each level}

A3 produces a consistent sequence of estimators for the population
bisimulation distances at each level:
\[
\widehat s_g \xrightarrow{|g|\to\infty} \mu_g, \qquad
\widehat s_{dom} \xrightarrow{|dom|\to\infty} \mu_{dom}.
\]
A2 produces estimators whose magnitude diverges with sample size:
$\widehat s_g^{\text{A2}} = |g|\cdot\widehat s_g^{\text{A3}} \to \infty$ as
$|g|\to\infty$ (assuming $\mu_g\neq 0$). For Variant 2 to admit a
well-defined population limit under A2, one must also let the
hyperparameters diverge in lockstep with the level size, which is the
formal statement of ``the scale is absorbed into $\lambda$''; it is a
\emph{rescaling of the mathematical objective} by a batch-composition
factor, not a consistent statistical model.

% =============================================================================
\section{The decoupling property in Context II, precisely}

Returning to the user's concern: if the subtraction $s_\tau-s_g$ (or
$s_g-s_{dom}$) appears in a context \emph{outside} the $\Ehat\cdot\hat
A_\tau$ form, we want it to carry information about the deviation of one
level from the next, independent of batch composition.

\begin{proposition}[Decoupling / interpretability of the subtraction]
\label{prop:decouple}
Under A3, $s_\tau - \widehat s_g$ has the interpretation of a
within-group deviation, with
\[
\Ehat_{\tau\in g}[s_\tau - \widehat s_g] = 0
\quad\text{and}\quad
\Ehat_{\tau\in g}[(s_\tau - \widehat s_g)^2] = \text{within-group sample variance of }s_\tau.
\]
Under A2, $s_\tau - \widehat s_g^{\text{A2}} = s_\tau - |g|\cdot\widehat s_g^{\text{A3}}$,
which has no interpretation as a deviation and is dominated by the
$|g|\cdot\widehat s_g^{\text{A3}}$ term for any $|g|\geq 2$.
\end{proposition}

\begin{corollary}[Penalty / novel-variant safety]
Any objective term of the form $\rho(s_\tau - \widehat s_g)$ with
$\rho:\mathbb R\to\mathbb R_{\geq 0}$ convex and $\rho(0)=0$ penalises
``trajectory drifts away from group centre'' under A3. Under A2, the
same functional form penalises approximately $\rho(|g|\cdot\widehat s_g)$
---a penalty on the group mean scaled by $|g|$, which has no relationship
to the intended per-trajectory deviation.
\end{corollary}

\begin{proof}
Direct consequence of Proposition~\ref{prop:decouple}.
\end{proof}

% =============================================================================
\section{Conclusion}

The user's diagnosis is correct: once we drop the GRPO zero-group-mean
assumption and allow the cross-hierarchy subtraction to appear in
contexts other than $\Ehat[\cdot]\hat A_\tau$, there \emph{is} a
well-defined and principled distinction between A2 and A3. A3 is not
merely preferable on engineering grounds (principle 2 / verl-alignment /
single-$\delta$) but \emph{uniquely} canonical in the precise sense of
Theorem~\ref{thm:unique}:

\begin{enumerate}
\item A3 is the only candidate that makes $\widehat s_g$ and
$\widehat s_{dom}$ empirical conditional means of their respective
children (C2, C3).
\item A3 is the only candidate that yields an ANOVA decomposition of
$s_\tau$ into within-group, between-group-within-domain, and domain
components (Proposition~\ref{prop:anova}).
\item A3 is the only candidate whose estimators converge to a
batch-composition-independent population limit (\S4.3).
\item Under A3 alone, a penalty of the form $\rho(s_\tau-\widehat s_g)$ or
$\rho(\widehat s_g-\widehat s_{dom})$ with $\rho$ convex, $\rho(0)=0$
admits the intended "deviation from level centre" interpretation
(Proposition~\ref{prop:decouple}).
\end{enumerate}

A2 passes the narrower Context-I recovery test
(Proposition~\ref{prop:ctx1}) but fails all four of the above
structural properties the moment the subtraction is used anywhere
outside the $\Ehat\cdot\hat A_\tau$ linear form. Absorbing the scale
into $\lambda$ is a valid rescaling but not a canonical choice: it makes
the hyperparameter a function of $|g|$ and $|dom|$, which is precisely
what ``canonical'' rules out.

The recommendation of A3 in \texttt{bisim\_eqs\_multidomain.tex}~\S3
therefore stands on two independent and additive foundations: the
engineering/infra alignment argument
(\texttt{bisim\_eqs\_recovery.tex}~\S7) and the canonical-decomposition
argument here.

\end{document}
